\scanpage{7}
Supposons $(C',D',\varphi')=(C,D,\varphi)$ donc $G=G'$, $I=I'$, $h'_i=h_i$

donc $\underline{\mathrm{Hom}}(,)=\underline{\mathrm{End}}((C,D,\varphi)$ est une Gr-catégorie.

\[
\pi_0\big(\underline{\mathrm{End}}(C,D,\varphi)\big)\overset{?}{=}
\]

\uline{NB} $(u_1,v_1,\alpha_1)\sim(u_2,v_2,\alpha_2)$ \deleted{\unsure{ssi}} implique

a) $u_1=u_2\ (=u)$

b) $v_2$ conjugué de $v_1$ i.e. $\exists\mu\in G'$ tel que $v_2=\mathrm{int}(\mu)v_1$

\[
\begin{aligned}
\text{d'où}\quad v_2(h_i)&=\mathrm{int}(\mu)\,v_1(h_i)=\mu\,v_1(h_i)\mu^{-1}\\[-2pt]
\Vert\quad &\hspace{11em}\Vert\\[-2pt]
\alpha_2(i)h_{u(i)}^{\deleted{\unsure{\prime}}}\alpha_2(i)^{-1}
&=\mu\alpha_1(i)h_{\unsure{u_1(i)|u_2(i)|u(i)}}^{\deleted{\unsure{\prime}}}\alpha_1(i)^{-1}\mu^{-1}
\end{aligned}
\]
d'où \quad $(\mathrm{int}\,\alpha_2(i))h_{u(i)}=(\mathrm{int}(\mu\alpha_1(i)))\cdot h_{u(i)}$

i.e \deleted{int} $\alpha_2(i)^{-1}\mu\alpha_1(i)^{\deleted{\unsure{-1}}}$ centralise $h_{u(i)}$

Or sauf le cas de la ``sphère à 1 seul trou'', chaque $h_j$ est un élément libre, et \inserted{en tous cas,} son centralisateur est le groupe engendré par $h_j$ ; \deleted{donc, dans les cas, \unsure{il}} donc a) et b) impliquent déjà l'existence des $\lambda_i$ tels que $\alpha_2(i)^{-1}\mu\alpha_1(i)=h_{u(i)}^{\lambda_i}$ (d'ailleurs uniques sauf le cas de la sphère \unsure{à} un seul trou i.e. $G=1$, $\mathrm{card}\ I=1$), donc
\[
(u_1,v_1,\alpha_1)=(u_2,v_2,\alpha_2)\quad\text{sss}\quad u_1=u_2,\ \widetilde v_2=\widetilde v_1
\]
($\sim$ dénote passage au groupe extérieur). Donc

$\pi_0(\underline{\mathrm{End}})\simeq$ groupe des couples $(u,\widetilde v)$, où $u\in\unsure{\mathfrak S_I|\mathfrak G_I}$, $\widetilde v\in\mathrm{Aut\,ext}\,G$, tels que $\forall i\in I$, on ait $\widetilde v(h_i)$ conjugué à $h_{u(i)}$ (pour $v$ un \unsure{auto|aut.} de $G$ qui représente $\widetilde v$).

\uline{NB} $\sigma$ est déjà déterminé par $\widetilde v$, ce qui résulte du fait que si $i,j\in I$, $i\ne j$ alors $h_i$ et $h_j$ sont non conjugués (en fait leurs images dans $G_{\rm ab}$ sont distinctes !) Donc $\pi_0(\underline{\mathrm{End}})\subset\mathrm{Aut\,ext}(G)$, est formé des $\widetilde u\in\mathrm{Aut\,ext}$ tels que $\forall i\in I$, $\widetilde v(h_i)$ est conjugué à un des $h_j$.
% archive sheet number: 6
